% Reproduced from the pinned incoming guarded_source report.
% Historical references to a proposed or new calculus refer to that report.
\section{The complete proposed presentation}
\label{sec:calculus}

\subsection{Terms, public formulas, and auxiliary formulas}

There is one numerical sort. Terms are finite expressions
\[
 t ::= x\mid 0\mid S(t)\mid I(n,a;u.h).
\]
Only $u$ is bound in $h$; it is not bound in $n$ or $a$. The body $h$ may
contain fixed numerical parameters. Terms do not contain a numerical
variable ranging over program descriptions, and no uniform evaluator is
postulated. Write $1=S0$, $2=S1$, and
\[
 C(v,a,b):=I(v,a;u.b),
\]
where the update variable is unused. All binders use capture-free
substitution and may be renamed.

\begin{definition}[Public language]
A canonical public formula is a literal equation $s=t$, or a formula
$D_{x,y}B$ satisfying one of the following two clauses:
\begin{enumerate}
\item $B$ is a compound canonical formula, so its root is $D$;
\item $B$ is exactly of the guarded atomic shape
\[
 B(x,y)\equiv I(y,0;u.h)=I(y,0;u.k).
 \tag{G}\label{eq:guard}
\]
\end{enumerate}
The bodies $h,k$ may depend on $x,y,u$ and passive parameters. In particular,
the guard is a syntactic condition, not a test for the theoremhood of a base
case. Both sides in \eqref{eq:guard} compute to zero at $y=0$.
The designated contradiction is the public equation $0=1$.
\end{definition}

The finite auxiliary formula language used \emph{inside proofs} is larger:
\[
 A ::= (s=t)\mid Q_yA\mid D_{x,y}A.
 \tag{A}\label{eq:auxiliary}
\]
It permits unguarded atomic matrices. A retained context is a finite reusable
list of auxiliary formulas. A public sequent has a canonical context and
canonical conclusion, but its proof may use \eqref{eq:auxiliary} internally.
This convention is part of the construction, not an omitted restriction.
There is no claim below about forcing every intermediate formula to satisfy
the public guard.

\subsection{All primitive axioms and rules}

In every context the two nonlogical axiom schemas are
\begin{align}
 I(0,a;u.h)&=a, \tag{I0}\label{eq:i0}\\
 I(Sn,a;u.h)&=h[I(n,a;u.h)/u]. \tag{IS}\label{eq:is}
\end{align}
There are exactly the following four inference patterns besides those two
computation schemas.

\paragraph{Hypothesis.}
\[
 \frac{}{\Gamma\vdash A}\qquad(A\in\Gamma).
 \tag{Hyp}
\]

\paragraph{Elimination selected by the source.}
For any fixed auxiliary formula template $F(v)$, put
\[
 M(s,t;F):=
 \begin{cases}
 Q_vF(v),&s\equiv0,\\
 s=t,&s\synneq0.
 \end{cases}
\]
The rule is
\[
 \frac{\Gamma\vdash M(s,t;F)\qquad\Gamma\vdash F(s)}
      {\Gamma\vdash F(t)}.
 \tag{SE}\label{eq:se}
\]
Here $s\synneq0$ means that the source is not literally the one-node term
$0$. A variable or an iterator that evaluates to zero is still a permitted
equality source. The $s\equiv0$ branch is ordinary anchored $Q$ elimination;
the other branch is restricted equality transport.

\paragraph{Paired elimination.}
\[
 \frac{\Gamma\vdash D_{x,y}B(x,y)\qquad
       \Gamma,Q_yB(b,y)\vdash A}
      {\Gamma\vdash A}.
 \tag{DE}\label{eq:de}
\]
The eigenparameter $b$ is fresh for $\Gamma$, $A$, and the passive
parameters of $B$.

\paragraph{One deterministic paired introduction.}
Define $\eta_B(z)=z$ when $B$ is a literal equation and $\eta_B(z)=0$
otherwise. Then
\[
 \frac{\Gamma,B(t,\eta_B(z))\vdash B(t,Sz)}
      {\Gamma\vdash D_{x,y}B(x,y)}.
 \tag{DI}\label{eq:di}
\]
The eigenparameter $z$ is fresh for $\Gamma$, $t$, and the passive
parameters of $B$. The local assumption is discharged. There is no
certificate choice between cases and progress: the root of the fixed
matrix determines it. Atomic matrices use progress; compound matrices use
their zero instance.

Call this auxiliary proof calculus $\UU$, and call its consequence relation
on public canonical sequents the proposed presentation. There are no
additional primitive axioms or rules for reflexivity, symmetry, successor
injectivity, successor separation, numerical coverage, induction,
propositional connectives, or ordinary quantifiers. Weakening, assumption
cut, and fresh renaming are admissible by induction on derivations. Arbitrary
numerical substitution into a restricted proof is not assumed in the
bootstrap argument.

\section{Soundness and the purpose of the guard}
\label{sec:sound}

Interpret the terms in iPRA by their primitive recursive functions, one fixed
finite iterator description at a time. Equality is unchanged. Give auxiliary
formulas the following compositional reading, denoted by a superscript $*$:
\begin{align*}
 (s=t)^*&=(s=t),\\
 (Q_yB)^*&=B^*(0)\to\forall y\,B^*(y),\\
 (D_{x,y}B)^*&=
 \begin{cases}
  \exists x\,\forall y\,(B(x,y)\to B(x,Sy)),&B\text{ a literal equation},\\
  \exists x\,(B^*(x,0)\to\forall y\,B^*(x,y)),&B\text{ compound}.
 \end{cases}
\end{align*}
This is a specified syntax-sensitive translation. It is not an assertion
that a single ordinary connective has interchangeable meanings when its
matrix is replaced by a logically equivalent formula.

\begin{lemma}[Auxiliary soundness]
If $\Gamma\vdash_{\UU}A$, iPRA proves the corresponding consequence
from $\Gamma^*$ to $A^*$.
\end{lemma}
\begin{proof}
Computation and hypotheses are immediate. Numerical substitution preserves
the root constructor of a formula, so the reading commutes with numerical
substitution; the equality branch of SE is sound. Its zero branch is
anchored $Q$ elimination.

At an atomic DI, the premise proves each progressive implication, and
universal introduction followed by existential introduction proves the
displayed atomic reading. At compound DI, numerical coverage supplies the
missing zero-or-successor split, yielding the anchored reading.

At an atomic DE, a progressive equation predicate gives
$E(0)\to\forall y\,E(y)$ by contextual equation induction in iPRA.
This is exactly the exported $Q$ resource. Compound DE is ordinary
existential elimination. All eigenvariable conditions are the standard
ones. Induction on the finite derivation completes the proof.
\end{proof}

The public ordinary reading, denoted by $A^\circ$, instead reads every
public $D$ uniformly as
\[
 (D_{x,y}B)^\circ=\exists x\,(B^\circ(x,0)\to\forall y\,B^\circ(x,y)).
 \tag{O}\label{eq:ordinary}
\]

\begin{lemma}[Agreement on the public grammar]
For every canonical formula $A$, iPRA proves $A^*\leftrightarrow A^\circ$.
\end{lemma}
\begin{proof}
At a guarded atomic matrix $E(x,y)$, computation proves $E(x,0)$.
Consequently, with $x$ fixed,
\[
 \forall y\,(E(x,y)\to E(x,Sy))
 \quad\longleftrightarrow\quad
 E(x,0)\to\forall y\,E(x,y).
\]
The forward implication is equation induction; the reverse uses the
theorem $E(x,0)$ to obtain all instances. Existential quantification
preserves the equivalence. At a compound matrix the two readings have the
same form, so structural induction applies.
\end{proof}

The guard is essential to this argument. Without a true zero instance, a
progressive predicate and an anchored predicate need not agree. For example,
the equation asserting that the parity iterator is $1$ is false at zero,
true at one, and false at two. Its anchored closure is true, while its
progressive closure is false. We make no completeness or composite claim
for all unguarded auxiliary formulas.

\section{The equational dependency theorem}
\label{sec:equational}

This section isolates a smaller set of hypotheses for the repository's
literal-zero-source argument \cite{source}. The singleton test and the
zero-filter construction are inherited from that argument. The dependency
statement, and its use in the combined presentation, are the point here.

\begin{definition}[Equational fragment]
Let $\EE$ have only equation formulas, hypotheses, the two computation
schemas, equation-template transport
\[
 \frac{\Gamma\vdash s=t\qquad\Gamma\vdash F(s)}{\Gamma\vdash F(t)}
 \quad(s\synneq0,\ F\text{ a literal equation template}),
 \tag{Tr$^-$}\label{eq:trminus}
\]
and contextual equation induction
\[
 \frac{\Gamma\vdash E(0)\qquad\Gamma,E(z)\vdash E(Sz)}
      {\Gamma\vdash E(r)}.
 \tag{EI}\label{eq:ei}
\]
The eigenparameter $z$ is absent from $\FV(\Gamma)$, the target $r$, and
all passive parameters of $E$; these remain fixed during induction.
\end{definition}

\begin{theorem}[Equational source elimination]
\label{thm:equational}
Unrestricted equation-template transport is admissible in $\EE$. Thus the
restriction $s\synneq0$ changes neither its closed equational theorems nor
its contextual equational consequences. No logical binder or logical case
introduction is needed.
\end{theorem}

\subsection{Equation cases and safe equality}

Equation cases already follow from EI: given $E(0)$ and $E(Sz)$ in a fixed
context, weaken the successor proof by the unused assumption $E(z)$ and
apply EI at the desired target. This observation removes the logical
case-introduction dependency from the original proof.

Reflexivity of an arbitrary $r$ follows by taking $b=I(0,r;u.0)$ and
using $b=r$ as both transport premises with motive $v=r$. The source $b$
is allowed. From $p=q$ with $p\synneq0$, permitted transport gives both
$q=p$ and $H(p)=H(q)$ for every numerical template $H$. All computation
equations have an iterator on the left and can therefore be reversed as
equations. Their reversed equations can be used as transport majors only
when the new left side is syntactically nonzero.

\begin{lemma}[Joining through a represented zero]
If $p,q\synneq0$ and $\Gamma$ proves $p=0$ and $q=0$, then $\Gamma$
proves $p=q$ using computation and restricted transport alone.
\end{lemma}
\begin{proof}
Put $z_0=I(0,0;u.0)$ and $b_p=I(0,p;u.0)$. Transport $b_p=b_p$
along $p=0$ in the left-hand occurrence of its initializer, obtaining
$z_0=b_p$. Reverse at the permitted source $z_0$, then use $b_p=p$
to get $p=z_0$. Similarly derive $q=z_0$. Reverse the latter at source
$q$ and transport through $z_0$ to obtain $p=q$. Every source used is
one of $p,q,z_0,b_p,b_q$, hence is permitted. This avoids transitivity
through literal zero before that operation has been established.
\end{proof}

\subsection{Two unconditional tests}

Define the following finite scalar terms:
\[
\begin{aligned}
 Z(v)&=C(v,1,0),& K(v)&=C(v,0,1),& H(v)&=C(v,1,2),\\
 N(n)&=I(n,0;u.Z(u)),& R(n)&=I(n,0;u.H(u)),& L(n)&=N(R(n)).
\end{aligned}
\]
These are explicit abbreviations over $0,S,I$, not new function symbols.

\begin{lemma}[Test and filter identities]
The equational fragment proves, at every term target $v$,
\[
 L(1)=1,\qquad L(SSv)=0,\qquad L(Sv)=Z(v),
 \tag{T}\label{eq:tests}
\]
and, at every term target $t$,
\[
 t=C(Z(t),t,0).
 \tag{F}\label{eq:filter}
\]
\end{lemma}
\begin{proof}
Equation cases prove $H(v)=S K(v)$: the zero endpoints compute to $1$
and the successor endpoints compute to $2$. These joins use the nonzero
numerals $1,2$. Computation gives $N(1)=1$, $N(2)=0$, and $R(1)=1$.
Also
\[
 R(Sv)=H(R(v))=S K(R(v)).
\]
Replacing the iterator source $R(Sv)$ in $R(SSv)=H(R(Sv))$ now gives
$R(SSv)=2$. Hence $L(1)=1$ and $L(SSv)=0$ by replacement at iterator
sources inside $N$.

For $L(Sv)=Z(v)$, use equation cases. At zero both endpoints equal $1$.
At a successor both equal zero, and both left-hand expressions are
iterator terms, so the preceding joining lemma applies. The case rule
itself supplies every arbitrary target; no substitution into a proof with
a variable transport source is required.

For the filter, apply equation cases to $v=C(Z(v),v,0)$. At zero the
right side computes to $C(1,0,0)=0$; at $Sv$ it computes to
$C(0,Sv,0)=Sv$. Reverse these equations at their iterator sources to
obtain the required branch orientations. Both branch contexts are empty.
\end{proof}

\subsection{Reversing an equality whose source is zero}

\begin{lemma}[Zero-source symmetry]
For every term $t$, $0=t\vdash_{\EE}t=0$.
\end{lemma}
\begin{proof}
If $t\equiv0$ syntactically, use reflexivity. Otherwise fix $t$ and retain the
assumption $0=t$. Put $A(n)=I(n,t;u.Su)$ and induct on
\[
 E(n)\equiv n=A(n).
\]
At the base, computation gives $A(0)=t$, and safe symmetry gives
$t=A(0)$. Transport in the retained $0=t$ at the permitted source $t$
gives $0=A(0)$.

At the step, $z=A(z)$ has a variable source, which is permitted.
Successor congruence gives $Sz=S A(z)$. Reverse the computation
$A(Sz)=S A(z)$ and transport at source $S A(z)$ to obtain $Sz=A(Sz)$.
The parameter $t$ and the context remain fixed.

Apply EI directly at $1$. It gives $1=A(1)$; forward computation of
$A(1)$ gives $1=St$. Congruence for $L$, followed by
\eqref{eq:tests}, gives $1=Z(t)$; safe symmetry at source $1$ gives
$Z(t)=1$. Weaken the unconditional filter \eqref{eq:filter} into the
retained context. Replace $Z(t)$ by $1$ there, then compute
$C(1,t,0)=0$. These last sources are iterators, and the result is $t=0$.
\end{proof}

\begin{proof}[Completion of Theorem~\ref{thm:equational}]
Full symmetry now follows by the safe-source case and the preceding lemma.
For congruence from $0=t$ with $t\synneq0$, first reverse it, use
permitted congruence at $t$, and reverse the resulting equality. The
literal-zero/literal-zero case uses reflexivity.

For transitivity $r=s$, $s=t$, a nonzero middle source is already allowed.
If $s$ is literal zero and one endpoint is literal zero, a premise gives
the conclusion. Otherwise reverse $0=t$ and join $r=0$, $t=0$ by the
represented-zero lemma. Thus all numerical equality algebra is available.
For $F(v)\equiv A(v)=B(v)$, use numerical congruence on both sides and
join with $A(s)=B(s)$ to obtain $A(t)=B(t)$.
\end{proof}

The proof uses weakening, cut, renaming, and only source-preserving helper
substitutions. The one fresh variable used as a transport source in the
induction step is never replaced by literal zero inside that derivation.
Arbitrary targets are supplied by the induction or case rule itself. This
distinction is necessary for a noncircular proof.

\section{Completeness of the combined presentation}
\label{sec:complete}

\begin{lemma}[Equation induction inside $\UU$]
Contextual EI is derivable in $\UU$ without using transport.
\end{lemma}
\begin{proof}
Given a progressive step for $E(y)$, take an unused witness coordinate
$x$. Atomic DI with witness $0$ gives $D_{x,y}E(y)$. DE opens
$Q_yE(y)$ independently of its fresh witness, and the zero-source branch
of SE applies it to the supplied base $E(0)$ at the desired target.
\end{proof}

This lemma may use an \emph{unguarded auxiliary} matrix $E$. It is why the
proof language was stated explicitly. Applying
Theorem~\ref{thm:equational}, with additional retained auxiliary assumptions
treated as passive context, restores all equation-template transport in
$\UU$ before any logical padding is needed.

Let $\GG$ denote the inherited guarded-core calculus: the same public
guarded formulas, the same deterministic introduction, and unrestricted
transport on canonical templates, with anchored $Q$ resources opened by DE.
The repository supplies effective translations, both formula composites,
and a whole-domain bi-interpretation of $\GG$ with iPRA \cite{guarded}.

\begin{lemma}[Transport on every canonical template]
\label{lem:canonicaltransport}
For a canonical template $F(v)$, the unrestricted rule
$s=t,F(s)\vdash_{\UU}F(t)$ is admissible.
\end{lemma}
\begin{proof}
Induct externally on the formula template, uniformly in retained contexts.
The equation case was just established. Suppose
$F(v)=D_{x,y}B(v,x,y)$. Open $F(s)$ at a fresh witness $a$.

If $B$ is guarded atomic, computation and the restored equality algebra
prove $B(s,a,0)$. Its opened $Q$ resource therefore yields $B(s,a,Sz)$.
Equation transport gives $B(t,a,Sz)$. Apply the deterministic atomic
introduction to the target at witness $a$, ignoring its local assumption
$B(t,a,z)$, and close DE.

If $B$ is compound, introduce the target by the deterministic case rule.
Use the smaller-template induction hypothesis and $t=s$ to turn its local
base $B(t,a,0)$ into $B(s,a,0)$. Apply the source $Q$ resource at $Sz$,
and use the induction hypothesis forward to obtain $B(t,a,Sz)$. Close
case introduction and DE. All newly chosen parameters are fresh for the
retained equality and context. Passive substitutions preserve the guard.
\end{proof}

In the inherited clause presentation, the canonical formulas in $\Gamma$
are read as ordinary assumptions, represented by their constant clauses.

\begin{theorem}[Exact public consequence relation]
\label{thm:consequence}
For every finite context $\Gamma$ of canonical formulas and every canonical
$A$,
\[
 \Gamma\vdash_{\UU}A\quad\Longleftrightarrow\quad\Gamma\vdash_{\GG}A.
\]
\end{theorem}
\begin{proof}
A guarded-core proof translates to $\UU$: its computations,
hypotheses, introductions, and resource eliminations are available, while
each unrestricted canonical transport is replaced by
Lemma~\ref{lem:canonicaltransport}. Starting from public formula contexts,
the reachable clauses are constant clauses for ordinary assumptions and
anchored clauses opened by DE. The former correspond to hypotheses, and
each use of the latter corresponds to a $Q$ elimination. No arbitrary
additional clause context is being identified with a $Q$ context.

Conversely, auxiliary soundness translates a $\UU$ proof to iPRA.
Agreement of the two readings on the canonical context and conclusion
gives its ordinary guarded reading. The inherited faithful guarded/iPRA
translation, including its canonical formula composite, returns a guarded
proof of the original sequent. This is an effective translation of finite
proofs using the inherited translations.
\end{proof}

The converse does not identify all auxiliary formulas with their ordinary
anchored readings. Only the public context and conclusion use the agreement
lemma. Likewise, Lemma~\ref{lem:canonicaltransport} does not assert full
transport or full formula congruence for every unguarded auxiliary binder.

\section{The two interpretations and both composites}
\label{sec:biinterpretation}

The consequence theorem transfers an existing bi-interpretation, rather
than merely transferring consistency. Here are its components and the
exact facts needed from the inherited foundations.

\subsection{The arithmetic bridge}

The ordinary scalar theory $T_{\rm eq}$ is intuitionistic first-order logic
with equality over the iterator terms, the two computation equations,
coverage $\forall x(x=0\lor\exists y\,x=Sy)$, separation $0=1\to\bot$,
and contextual induction for literal equations. Its inherited construction
derives predecessor, equality reflection into a Boolean term,
quantifier-free induction, pairing with projections, and compilation of
every fixed primitive recursive definition, in that order
\cite{foundations,lean}.

The direction to iPRA interprets each fixed iterator body by primitive
recursion. In the reverse direction, a description of a primitive recursive
function is compiled to an iterator term; a counter and the current result
can be packed into the scalar state. The translations prove both term
round trips equal to the original terms. They are interpretations over the
whole numerical domain with actual equality, not a restriction to standard
numerals or a quotient by an external equivalence relation.

The pinned Lean modules \code{BiInterpretation},
\code{NativeBiInterpretation}, and their dependencies state the two
directions and the round trips for the ordinary scalar and native theories.
The guarded grammar and padding bridge remain written mathematics in the
pinned research package. This report does not claim that the present
composition has been checked by Lean.

\subsection{The logical bridge and guarded padding}

The inherited native core has equations, $Q$, and $D$. Write
$\operatorname{sel}(v,A,B)$ for its formula selector, equivalent to $A$
at $v=0$ and to $B$ at successors, and $\operatorname{gate}(v,A)$ for its
formula gate, true at zero and equivalent to $A$ at successors. Their
equational leaves use the scalar selector $C$; syntactic padding aligns
compound shapes. The ordinary logical operations are macros, for example
\[
\begin{aligned}
 \bot&:=0=1,\\
 A\to B&:=Q_y\operatorname{sel}(y,A,B),\\
 A\wedge B&:=Q_y\operatorname{gate}
       (y,\operatorname{sel}(\operatorname{par}(y),A,B)),\\
 \exists x\,A(x)&:=D_{x,y}\operatorname{gate}(y,A(x)),\\
 A\vee B&:=\exists v\,\operatorname{sel}(v,A,B),\\
 \forall x\,A(x)&:=A(0)\wedge Q_yA(y),
\end{aligned}
\]
where the parity iterator is a fixed scalar term. All natural deduction
rules for these macros are inherited derived rules; they are not additions
to Section~\ref{sec:calculus}.

To pass explicitly to the equivalent one-object-binder presentation,
recursively replace each $Q_yB$ by $D_{u,y}B$ with a fresh unused $u$.
The two are interderivable in the old core: paired elimination gives $Q$
from the dummy-witness $D$, while case introduction and QE give the reverse
direction with witness $0$. After this expansion, $Q$ is needed only as a
context resource. For an equation $E$, put
\[
 G_vE=(C(v,0,s)=C(v,0,t))\quad(E\equiv s=t),\qquad
 L(E)=D_{u,v}G_vE,
\]
with fresh $u,v$, and let $P(E)=L(E)$, $P(A)=A$ at compound formulas.
Guarded padding gives $P(A)\dashv\vdash A$. The translation from the
unrestricted one-binder core to the guarded grammar is
\[
 T(s=t)=s=t,\qquad T(D_{x,y}B)=D_{x,y}P(TB).
\]
It commutes with numerical substitution up to renaming. The inherited
guarded-core argument proves both $JT(A)\dashv\vdash A$ in the old core
and $TJ(B)\dashv\vdash B$ for canonical guarded $B$, where $J$ is inclusion.
The guarded atomic case of the second composite uses its true base; it is
not an unrestricted congruence assertion for all auxiliary matrices.

\begin{theorem}[Bi-interpretability and equiconsistency]
\label{thm:biinterp}
The proposed public presentation, realized by \eqref{eq:ordinary}, is
bi-interpretable with full first-order iPRA. Both numerical domains are
whole and both numerical composite isomorphisms are given by $x=y$.
Both formula composites are provably equivalent to the original formulas.
In particular,
\[
 \varnothing\vdash_{\UU}0=1
 \quad\Longleftrightarrow\quad
 \IPRA\vdash\bot.
\]
\end{theorem}
\begin{proof}
Let $\Phi$ be the inherited iPRA-to-guarded translation, obtained by
composing the scalar compiler, the native logical translation, the
dummy-witness expansion of $Q$, and guarded padding. Let $\Psi$ be the
guarded ordinary reading followed by the
scalar-to-iPRA interpretation. Use these same maps for the new presentation.
Theorem~\ref{thm:consequence} preserves exactly every canonical consequence
used by $\Phi$, $\Psi$, and their guarded composite proofs.

The inherited iPRA composite proves
$\varphi\leftrightarrow\Psi\Phi(\varphi)$ for every ordinary formula
$\varphi$. The inherited guarded composite proves
$A\dashv\vdash\Phi\Psi(A)$ for every canonical $A$; the consequence
theorem makes this a proof in the proposed presentation. On numerical
domains nothing has changed, and inherited term round-trip equations
witness the literal identity graph. Finally the designated contradiction
has the same reading throughout, giving the displayed equivalence.
\end{proof}

This proof uses the inherited full bridge as a proved input, with its
sources supplied for inspection. A new derivation of predecessor or a new
primitive-recursion compiler is neither needed nor claimed.

